% Introduction rewritten after studying the five user-designated IEEE GRSL examples.
% Citation keys are provisional and must be synchronized with the final .bib file.

\section{Introduction}

The measurement and simulation of radar backscattering from ocean surfaces are important for maritime remote sensing, radar echo simulation, and target detection. Sea-surface roughness is multiscale: long gravity waves determine the large-scale tilt and motion of the surface, whereas shorter wind waves modify local slopes and the radar-resonant components. Consequently, the observed return depends jointly on sea state, incidence angle, radar wavelength, and surface geometry. High-resolution measurements further show that sea clutter can exhibit pronounced non-Gaussian tails and intermittent sea spikes, which may be confused with target responses \cite{carretero2010statistical,fiche2015xband,bocquet2022lga,melief2006seaspikes,rosenberg2013seaspike}. Physically consistent sea-surface realizations are therefore required to generate controllable radar echoes and to clarify the geometric origin of these anomalous returns.

Directional wave spectra are widely used to describe the multiscale energy distribution of wind-driven seas. The Elfouhaily spectrum provides a unified representation of gravity and gravity--capillary waves and has been adopted in many sea-surface scattering studies \cite{elfouhaily1997unified}. However, a linear spectral realization is approximately Gaussian and mainly preserves second-order statistics. Nonlinear hydrodynamic interactions sharpen wave crests, smooth wave troughs, and alter the distributions of slopes and curvatures, thereby affecting the scattering coefficient and Doppler spectrum \cite{zhang2023secondorder}. Higher-order sea-surface statistics are also important because the skewness function and its small-scale-wave parameter govern the upwind--downwind asymmetry of radar backscattering \cite{xie2024skewness}. Moreover, short-wave spectral components and their convolutions can contribute appreciably to the normalized backscattering cross section under different radar frequencies, wind speeds, and incidence angles \cite{xie2022convolutions}. These results indicate that a radar-oriented sea model should represent not only the total spectral energy but also nonlinear crest geometry and the directional organization of short waves.

The scattering from ordinary rough surfaces is commonly modeled by quasi-specular, Bragg, or two-scale formulations. Their relative contributions vary with incidence angle, and both mechanisms must be considered in the transition range \cite{ye2022improvedTSM}. Radar observations of marine film slicks also suggest that the source of short wind waves can be direction dependent and cannot always be described by conventional wind-input terms alone \cite{sergievskaya2024shortwaves}. For steep and breaking waves, additional mechanisms become important. Numerical and experimental studies have related low-grazing-angle sea spikes to steepening or breaking crests \cite{holliday1998steepening,sletten2003breaking,west2002plunging}, while full-wave and hybrid models have investigated crest diffraction, rough-face scattering, and multipath interactions \cite{west1998mmgtd,west1999ray,west2002multipath,zhao2005spilling,yang2011breaking}. Facet-based models have further introduced breaking-wave contributions into large sea surfaces and analyzed their effects on NRCS, sea spikes, and synthetic aperture radar images \cite{li2017facet}. A three-dimensional plunging breaker combined with ray tracing was also shown to generate target-like peaks in radar cross section and high-resolution range profiles \cite{dong2023plunging}.

Although these studies demonstrate the electromagnetic importance of nonlinear and breaking waves, most of them emphasize scattering from prescribed breaker profiles or numerically generated local surfaces. A computationally efficient procedure that connects a statistically constrained nonlinear background, directionally organized short wind waves, and an explicit localized curling crest is still needed. Directly adding an independent short-wave field may duplicate spectral energy, whereas a single-valued nonlinear height field cannot represent the forward extension and overturning facets of a plunging crest. These limitations make it difficult to preserve large-area sea statistics and simultaneously reproduce the local geometry responsible for strong radar returns.

This letter proposes a hierarchical sea-surface model for radar echo simulation. First, an Elfouhaily directional spectrum and a wind-dependent nonlinear Lie transformation are used to construct the background surface. Second, a selected short-wave band is replaced by an equal-energy directional component to form wind-wave ridges with a prescribed dominant propagation direction. Third, steep candidate crests are identified from elevation, propagation-direction slope, and curvature, and a compactly supported coordinate transformation is applied to generate a localized forward-extending and downward-curling breaker. The background and short-wave-enriched surfaces are evaluated using spectral consistency and along-wind and crosswind mean-square slopes. The local breaker is assessed using crest geometry, facet-orientation statistics, radar-facing projected area, and paired unit-footprint scattering enhancement. Finally, radar echoes generated from the complete surface are examined through range--time intensity, temporal correlation, and amplitude-distribution statistics. Section II presents the hierarchical surface model. Section III describes the validation metrics and radar echo generation. Section IV reports the numerical and experimental comparisons, and Section V concludes this letter.
